\begin{table}[!htbp]
\centering
\caption{So sánh bản chất và vai trò của các thành phần mật mã trong PKI}
\label{tab:crypto_comparison}
\small
\renewcommand{\arraystretch}{1.3}
\begin{tabularx}{\textwidth}{@{}l >{\raggedright\arraybackslash}p{3.2cm} >{\raggedright\arraybackslash}p{2.6cm} >{\raggedright\arraybackslash}X >{\raggedright\arraybackslash}p{3.2cm}@{}}
\toprule
\textbf{Khái niệm} & \textbf{Bản chất kỹ thuật} & \textbf{Khóa bí mật ($SK$)} & \textbf{Giá trị an ninh cung cấp} & \textbf{Hạn chế nội tại} \\
\midrule
\textbf{Cặp khóa ($PK, SK$)} & 
Hai khóa có liên hệ toán học; $PK$ công khai, $SK$ bảo mật. & 
Chỉ duy nhất chủ sở hữu hợp pháp nắm giữ và bảo vệ. & 
Nền tảng cho ký số, xác thực sở hữu và thỏa thuận khóa phiên. & 
Không tự liên kết được với danh tính thực thể ngoài đời thực. \\
\midrule
\textbf{Hàm băm ($H$)} & 
Ánh xạ thông điệp tùy ý thành digest có độ dài cố định. & 
Không sử dụng khóa. & 
Bảo đảm tính toàn vẹn, tóm lược dữ liệu, kháng va chạm. & 
Không xác thực nguồn gốc do bất kỳ ai cũng tính toán được. \\
\midrule
\textbf{Chữ ký số} & 
Cặp thuật toán Sign/Verify thao tác trên giá trị băm thông điệp. & 
Ký bằng $SK$ của chủ thể; kiểm tra bằng $PK$ tương ứng. & 
Xác thực nguồn gốc dữ liệu, bảo đảm tính toàn vẹn và chống chối bỏ. & 
Chỉ xác minh tính hợp lệ của dữ liệu với $PK$, không định danh người ký. \\
\midrule
\textbf{Chứng thư số} & 
Cấu trúc dữ liệu liên kết $PK$ với danh tính, được CA ký. & 
Ký bởi $SK_{CA}$; cấu trúc chứng thư không bao giờ chứa $SK$. & 
Thiết lập liên kết tin cậy giữa khóa công khai và danh tính chủ thể. & 
Phụ thuộc độ tin cậy của CA, thời hạn hiệu lực và trạng thái thu hồi. \\
\bottomrule
\end{tabularx}
\end{table}
